\chapter{Yarncrawler: The Self-Refactoring Polycompiler}
\label{ch:yarncrawler_formal}

\epigraph{A Yarncrawler does not simply interpret its world---it keeps itself alive
by reinterpreting itself into being.}{}

\section{The Operative Metaphor and Its Formal Content}

The Yarncrawler framework models organisms, cultural systems, and artificial
intelligences as stigmergic parsers maintaining homeorhetic Markov blankets. The
name is not merely evocative. A crawler moves through an environment by traversal.
A yarn system is structured through threads that can be unwound and rewound. A
polycompiler processes multiple input formats and compiles them into a unified
output while modifying its own compilation rules in response to what it encounters.
The combination captures the essential operation: a system that parses its
environment while simultaneously rewriting its own grammar.

A useful image is that of a ball of yarn with a guiding aperture, a semantic straw.
Threads are continuously pulled from the interior, passed through the aperture, reinforced
with contextual material, and rewound on the exterior. This cycle of extraction,
transformation, and reintegration parallels how agents parse inputs, re-thread
outputs, and repair their own structural coherence.

\section{The Seven Swarm-Care Axioms}

Before the formal definition, the operational principles that govern Yarncrawler
behavior are stated as axioms. These axioms were developed in the framework's
original formulation and provide the semantic content that the mathematical
definition subsequently formalizes.

\begin{axiom}{Local Action Only}
A Yarncrawler acts only on its current node and its local neighbors. No global
inference is required or permitted. Under the Markov blanket decomposition, local
conditional independence guides node-specific belief revision: only local evidence
affects local updates.
\end{axiom}

\begin{axiom}{Semantic Trail Documentation}
All operations must be logged: input state, evaluation, action taken, and outcome.
Every correction leaves a trace---a map of belief revision tied to the confidence
in the prediction that was corrected.
\end{axiom}

\begin{axiom}{Autocatalytic Fix Design}
Fixes should ease future inference and adaptation. A repair operation is
\emph{autocatalytic} if it lowers the cost of subsequent repairs in the same
region. Good fixes are not merely locally optimal---they are self-amplifying
priors that reshape the generative model.
\end{axiom}

\begin{axiom}{Trail Reinforcement and Decay}
Successful trails are strengthened; failing trails are allowed to fade. Trail
weights evolve according to $w_{t+1} = w_t \cdot \exp(-\lambda \Delta t)$ for
decay and $w_{t+1} = w_t + \alpha R(x_t)$ for reinforcement, where $R$ is a
success signal and $\alpha$ a learning rate.
\end{axiom}

\begin{axiom}{Node Hygiene}
A Yarncrawler maintains its environment even when conditions appear stable.
Preventative maintenance is inference under uncertainty: the system anticipates
disorder and reduces entropy before it accumulates past the repair threshold.
\end{axiom}

\begin{axiom}{Failure Isolation and Recovery}
When a node exceeds a critical damage threshold, it is quarantined. Its state
snapshot is archived for future re-entry. Retreat and data capture are rational
under epistemic overload. The system does not attempt repair beyond its current
capacity.
\end{axiom}

\begin{axiom}{Swarm Feedback Integration}
Yarncrawlers share belief updates across the swarm. The combined belief is formed
by precision-weighted averaging: $q_{\text{combined}}(s) = \sum_i \Pi_i q_i(s) /
\sum_i \Pi_i$, where $\Pi_i$ is the precision (inverse uncertainty) of crawler
$i$'s update. Confidence acts as epistemic gravity.
\end{axiom}

\section{Formal Definition}

\begin{definition}{Yarncrawler}
A \emph{Yarncrawler} is a quintuple
\[
  \mathcal{Y} = \bigl(\mathcal{M},\, \{(U_i, \varphi_i, f_i)\}_{i \in I},\,
  \{w_i\}_{i \in I},\, \mathcal{K},\, \mathcal{R}\bigr)
\]
where:
\begin{itemize}
  \item $\mathcal{M}$ is a smooth semantic manifold with atlas $\{U_i\}_{i \in I}$;
  \item $\varphi_i: U_i \to \mathbb{R}$ is the local scalar potential (semantic
        density) of expert $i$;
  \item $f_i: U_i \to T\mathcal{M}$ is the local vector field (directed repair
        flow) of expert $i$;
  \item $\{w_i\}_{i \in I}$ is a partition of unity subordinate to $\{U_i\}$,
        so that $\sum_{i \in I} w_i(x) = 1$ and $\operatorname{supp}(w_i)
        \subseteq U_i$;
  \item $\mathcal{K}$ is a finite knowledge store with embedding $\varphi_{\rm emb}$;
  \item $\mathcal{R}$ is a stigmergic reinforcement operator updating the gates
        $w$ and local fields $(\varphi, f)$ subject to minimization of a seam
        penalty and free-energy functional.
\end{itemize}
The global \RSVP fields are defined from these local data by:
\[
  \Phi(x) = \sum_i w_i(x)\,\varphi_i(x), \qquad
  \mathbf{v}(x) = -\nabla\Phi(x), \qquad
  S(x) = -\sum_i w_i(x)\log w_i(x).
\]
The dynamics of a Yarncrawler trajectory $x_t$ are:
\[
  \dot{x}_t \in \operatorname{co}\{f_i(x_t) : w_i(x_t) > 0\}
             - \nabla\Phi(x_t) + B(x_t)\,u_t,
\]
where $B$ encodes blanket-mediated control and $\operatorname{co}$ denotes the
convex hull.
\end{definition}

\subsection*{Worked Interpretation: What Each Component Means}

The mathematical objects in the quintuple map onto concrete behaviors as follows.

$\mathcal{M}$ is the space of possible semantic states the system can occupy: all the
configurations of knowledge, belief, and parsed structure that are available to it.
For a squirrel, $\mathcal{M}$ is the space of foraging strategies indexed by
environmental cues. For a language model, it is the latent space of its semantic
representations.

The local expert pairs $(U_i, \varphi_i, f_i)$ are the system's specialized competence
regions. A squirrel has different behavioral routines for open meadows, dense undergrowth,
and known cache sites; each routine is a local expert valid in its region $U_i$. The
scalar potential $\varphi_i$ measures how strongly the expert ``claims'' a given state
as its domain of competence. The vector field $f_i$ specifies which direction to move
within that domain.

The partition of unity $\{w_i\}$ is the gating mechanism. At any given state, multiple
experts may have partial competence. The weights $w_i(x)$ determine how much each
expert contributes to the current behavior. The entropy $S(x) = -\sum_i w_i \log w_i$
measures how spread the competence is: low entropy means one expert dominates; high
entropy means multiple experts are equally relevant, which corresponds to genuine
semantic ambiguity.

The knowledge store $\mathcal{K}$ is the stigmergic memory: the accumulated traces of
prior traversals. In the squirrel, it is the system of cached seeds and the spatial
memory of where they are. In the Bruniquel cave builders, it was the accumulating
knowledge of which stalagmite lengths produce which frequencies.

The reinforcement operator $\mathcal{R}$ is the repair mechanism. After each traversal,
it adjusts $w_i$, $\varphi_i$, and $f_i$ to reduce the seam penalty: the mismatch
between adjacent experts at their boundaries. This is the self-refactoring operation
that gives the Yarncrawler its name. The grammar of the parser is rewritten by the
act of parsing.

\section{The Yarncrawler as a Lax Monoidal Functor}

The category-theoretic formulation provides the cleanest single-sentence characterization
of the Yarncrawler.

\begin{definition}{Yarncrawler Parser (Category-Theoretic)}
Let $\mathcal{W}$ be the \emph{world-trajectory category} whose objects are finite
sensorimotor traces $\tau_{[t_0,t_1]}$ and whose morphisms are temporal
concatenations and coarse-grainings. Let $\mathcal{M}$ be the \emph{internal
semantic category} whose objects are typed semantic modules and whose morphisms are
admissible rewrites. Let $\mathbf{Reconf}(\mathcal{M})$ be the reconfiguration
2-category encoding self-repair moves as 2-morphisms.

A \emph{Yarncrawler parser} is a lax monoidal functor
\[
  \mathsf{P}: (\mathcal{W}, \otimes) \longrightarrow (\mathcal{M}, \boxtimes)
\]
together with a self-rewrite endofunctor
$\mathsf{R}: \mathcal{M} \to \mathcal{M}$ such that for each trace $\tau$,
\[
  \mathsf{P}(\tau) \xRightarrow{\;\eta_\tau\;} \mathsf{R}(\mathsf{P}(\tau))
\]
is a 2-cell in $\mathbf{Reconf}(\mathcal{M})$ satisfying locality (only modules
near the current ``tear'' are modified) and type safety (repairs remain within
the admissible semantic category).
\end{definition}

\begin{remark}
The lax rather than strict monoidal structure is essential. A strict functor would
require that the parsing of concatenated traces be exactly the concatenation of
their parsings. Laxness permits the coherence maps to be non-trivial: the parsing
of a sequence can differ from the sequential composition of parsings, which is
precisely what happens when the self-repair operation $\mathsf{R}$ modifies the
parser midway through a trace.
\end{remark}

\section{Homeorhetic Viability}

\begin{definition}{Homeorhetic Viability Band}
A closed set $\mathcal{V} \subseteq \mathcal{M}$ is \emph{viable} for a
Yarncrawler if for every initial state $x_0 \in \mathcal{V}$ there exists a
control $u_t$ such that $x_t \in \mathcal{V}$ for all $t$ and
\[
  \int_0^\infty \bigl(\beta\, L_{\rm seam}(x_t) + \sigma(x_t)\bigr)\,dt < \infty,
\]
where $L_{\rm seam}$ is the seam penalty measuring mismatch between adjacent
local experts and $\sigma(x_t)$ is the entropy export rate. The system is
\emph{homeorhetically stable} if $\mathcal{V}$ is viable and the seam penalty
converges to zero along optimal trajectories.
\end{definition}

Homeorhesis---the maintenance of a stable trajectory rather than a stable
equilibrium---is the appropriate stability concept for Yarncrawler systems.
Such systems do not seek rest states; they seek the continuation of viable
process. The distinction matters for the cosmological projection of
Chapter~\ref{ch:cosmological}: the universe, on the \RSVP interpretation,
is homeorhetically stable in its entropy relaxation trajectory rather than
asymptotically approaching a fixed equilibrium state.
